\subsection{SYMMETRIC ALGEBRA OF A DIRECT SUM. SYMMETRIC ALGEBRA OF A FREE MODULE. SYMMETRIC ALGEBRA OF A GRADED MODULE}

Let A be a commutative ring, \(M = \bigoplus_{\lambda \in L} M_{\lambda}\) the direct sum of a family of A-modules and \(j_{\lambda}: M_{\lambda} \to M\) the canonical injection; we derive an A-homomorphism of algebras \(\mathsf{S}(j_{\lambda}): \mathsf{S}(\mathsf{M}_{\lambda}) \to \mathsf{S}(\mathsf{M})\) . Since \(\mathsf{S}(\mathsf{M})\) is commutative, Proposition 8 of §4, no. 5, can be applied to the homomorphisms \(\mathsf{S}(j_{\lambda})\) and there therefore exists a unique algebra homomorphism

\[
g \colon \bigotimes_ {\lambda \in L} \mathbf {S} (\mathbf {M} _ {\lambda}) \rightarrow \mathbf {S} (\mathbf {M})\tag{10}
\]

(also denoted by \(g_{\mathbf{M}}\) ) such that \(\mathbf{S}(j_{\lambda}) = g \circ f_{\lambda}\) for all \(\lambda \in \mathbf{L}\) , where

\[
f _ {\lambda} \colon \mathbf {S} (\mathbf {M} _ {\lambda}) \rightarrow \bigotimes_ {\lambda \in \mathbf {L}} \mathbf {S} (\mathbf {M} _ {\lambda})
\]

denotes the canonical homomorphism.

PROPOSITION 9. The canonical homomorphism \(g\) (formula (10)) is a graded algebra isomorphism (cf.~§4, no. 8, Remark 1).

To prove that \(g\) is bijective, we define an algebra homomorphism

\begin{enumerate}
\def\labelenumi{(\arabic{enumi})}
\setcounter{enumi}{10}
\tightlist
\item
\end{enumerate}

\[
h \colon \mathbf {S} (\mathbf {M}) \rightarrow \bigotimes_ {\lambda \in L} \mathbf {S} (\mathbf {M} _ {\lambda})
\]

such that \(g \circ h\) and \(h \circ g\) are respectively the identity mappings on \(\mathbf{S}(\mathbf{M})\) and \(\bigotimes_{\lambda \in \mathbf{L}} \mathbf{S}(\mathbf{M}_{\lambda})\) . For each \(\lambda \in \mathbf{L}\) , let \(u_{\lambda}\) be the composite linear mapping

\[
\mathbf {M} _ {\lambda} \xrightarrow {\phi_ {\mathrm{M} _ {\lambda}} ^ {\prime}} \mathbf {S} (\mathbf {M} _ {\lambda}) \xrightarrow {f _ {\lambda}} \bigotimes_ {\lambda \in L} \mathbf {S} (\mathbf {M} _ {\lambda}).
\]

There exists one and only one A-linear mapping \(u: \mathbf{M} \to \bigotimes_{\lambda \in L} \mathbf{S}(\mathbf{M}_{\lambda})\) such that \(u \circ j_{\lambda} = u_{\lambda}\) for all \(\lambda \in L\) . As the \(\mathbf{S}(\mathbf{M}_{\lambda})\) are commutative, so is their tensor product (§ 4, no. 5) and hence (no. 1, Proposition 2) there exists a unique algebra homomorphism \(h: \mathbf{S}(\mathbf{M}) \to \bigotimes_{\lambda \in L} \mathbf{S}(\mathbf{M}_{\lambda})\) such that \(h \circ \phi_{\mathbf{M}}' = u\) ; on the other hand, it is immediate that \(u(\mathbf{M})\) is contained in the submodule of elements of degree 1 of the graded algebra \(\bigotimes_{\lambda \in L} \mathbf{S}(\mathbf{M}_{\lambda})\) and hence \(h\) is a graded algebra homomorphism. For \(x_{\lambda} \in \mathbf{M}_{\lambda}\) ,

\[
h \left(g \left(u _ {\lambda} \left(x _ {\lambda}\right)\right)\right) = h \left(g \left(f _ {\lambda} \left(\phi_ {M _ {\lambda}} ^ {\prime} \left(x _ {\lambda}\right)\right)\right)\right) = h \left(S \left(j _ {\lambda}\right) \left(\phi_ {M _ {\lambda}} ^ {\prime} \left(x _ {\lambda}\right)\right)\right) = h \left(\phi_ {M} ^ {\prime} \left(j _ {\lambda} \left(x _ {\lambda}\right)\right)\right) = u _ {\lambda} \left(x _ {\lambda}\right);
\]

as the \(u_{\lambda}(x_{\lambda})\) generate the algebra \(\bigotimes_{\lambda \in L} S(M_{\lambda})\) (§ 4, no. 5, Proposition 8), \(h \circ g\) is certainly the identity mapping. Similarly,

\[
g \left(h \left(\phi_ {\mathrm{M}} ^ {\prime} \left(j _ {\lambda} \left(x _ {\lambda}\right)\right)\right)\right) = g \left(u _ {\lambda} \left(x _ {\lambda}\right)\right) = g \left(f _ {\lambda} \left(\phi_ {\mathrm{M} _ {\lambda}} ^ {\prime} \left(x _ {\lambda}\right)\right)\right) = \mathrm{S} \left(j _ {\lambda}\right) \left(\phi_ {\mathrm{M} _ {\lambda}} ^ {\prime} \left(x _ {\lambda}\right)\right) = \phi_ {\mathrm{M}} ^ {\prime} \left(j _ {\lambda} \left(x _ {\lambda}\right)\right)
\]

and, as the elements \(\phi_{\mathbf{M}}^{\prime}(j_{\lambda}(x_{\lambda}))\) generate the algebra \(\mathbf{S}(\mathbf{M})\) , \(g \circ h\) is certainly the identity mapping.

Remark (1) Let \(\mathbf{N} = \bigoplus_{\lambda \in L} \mathbf{N}_{\lambda}\) be the direct sum of another family of A-modules with \(L\) as indexing set and, for all \(\lambda \in L\) , let \(v_{\lambda}: \mathbf{M}_{\lambda} \to \mathbf{N}_{\lambda}\) be an A-linear mapping, whence there is an A-linear mapping \(v = \bigoplus_{\lambda} v_{\lambda}: \mathbf{M} \to \mathbf{N}\) (II, § 1, no. 6, Proposition 6). Then the diagram

\[
\begin{array}{c} \bigotimes_ {\lambda \in L} S (M _ {\lambda}) \xrightarrow {g _ {M}} S (M) \\ \bigoplus_ {\lambda \in L} S (v _ {\lambda}) \Biggl \downarrow \qquad \qquad \qquad \qquad \qquad \qquad \qquad \qquad \qquad \Biggl \downarrow S (v) \\ \bigotimes_ {\lambda \in L} S (N _ {\lambda}) \xrightarrow {g _ {N}} S (N) \end{array}
\]

is commutative, as follows from the definitions (§ 4, no. 5, Corollary to Proposition 8).

The sub-A-module of \(\bigotimes_{\lambda \in L} \mathbf{S}(\mathbf{M}_{\lambda})\) with which \(\mathbf{S}^n(\mathbf{M})\) is identified by means of the isomorphism \(g\) can be described more precisely. For every finite subset \(J\) of \(L\) , we write \(\mathbf{E}_J = \bigotimes_{\lambda \in J} \mathbf{S}(\mathbf{M}_{\lambda})\) , so that \(\bigotimes_{\lambda \in L} \mathbf{S}(\mathbf{M}_{\lambda}) = \varinjlim \mathbf{E}_J\) relative to the directed set \(\mathfrak{F}(\mathrm{L})\) of finite subsets of \(\mathbf{L}\) , by definition (§ 4, no. 5). For every family \(\nu = (n_{\lambda}) \in \mathbf{N}^{(\mathbf{L})}\) (thus having finite support) such that \(\sum_{\lambda \in \mathbf{L}} n_{\lambda} = n\) and every finite subset \(\mathbf{J}\) of \(\mathbf{L}\) containing the support of the family \(\nu\) , we write

\[
\mathbf {S} ^ {\mathrm{J}, \nu} (\mathbf {M}) = \bigotimes_ {\lambda \in J} \mathbf {S} ^ {n _ {\lambda}} (\mathbf {M} _ {\lambda})\tag{12}
\]

so that the submodule \(E_{J,n}\) of elements of degree n in \(E_{J}\) is the direct sum of the \(S^{J,v}(M)\) over all the families v of support contained in J and such that \(\sum_{\lambda\in L}n_{\lambda}=n\) ( \(\S4\) , no. 7, Proposition 10 and \(\S4\) , no. 8). As a convention we write \(S^{J,v}(M)=\{0\}\) for the families v whose support is not contained in J; then \(E_{J,n}\) can also be called the direct sum of all the \(S^{J,v}(M)\) , where v runs through the set \(H_{n}\) of all families \(v=(n_{\lambda})_{\lambda\in L}\) such that \(\sum_{\lambda\in L}n_{\lambda}=n\) . Since \(S^{0}(M_{\lambda})\) is identified with A, clearly also, for two finite subsets \(J\subset J'\) of L and a family v of support contained in J, the canonical mapping \(S^{J,v}(M)\to S^{J',v}(M)\) (restriction of the canonical mapping \(E_{J}\to E_{J'}\) to \(S^{J,v}(M)\) ) is bijective. If we write, for all \(v\in H_{n}\) ,

\[
\mathbf {S} ^ {\nu} (\mathbf {M}) = \varinjlim \mathbf {S} ^ {\mathrm{J}, \nu} (\mathbf {M})\tag{13}
\]

it is seen that, taking account of II, § 6, no. 2, Proposition 5:

COROLLARY. The A-module \(\mathbf{S}^n (\mathbf{M})\) is the image under isomorphism (10) of the submodule of \(\bigotimes_{\lambda \in \mathbf{L}}\mathbf{S}(\mathbf{M}_{\lambda})\) the direct sum of the submodules \(\mathbf{S}^{\nu}(\mathbf{M})\) for all the families \(\nu = (n_{\lambda})\in \mathbf{N}^{(L)}\) such that \(\sum_{\lambda \in \mathbf{L}}n_{\lambda} = n\) ; if \(\mathbf{J}\) is the support of \(\nu\) , \(\mathbf{S}^{\nu}(\mathbf{M})\) is canonically isomorphic to \(\bigotimes_{\lambda \in \mathbf{J}}\mathrm{S}^{n_{\lambda}}(\mathbf{M}_{\lambda})\) .

In general \(\mathbf{S}^{\nu}(\mathbf{M}),\bigotimes_{\lambda \in J}\mathbf{S}^{n_{\lambda}}(\mathbf{M}_{\lambda})\) and their image in \(\mathbf{S}^n (\mathbf{M})\) are identified.

THEOREM 1. Let A be a commutative ring and M a free A-module with basis \((e_{\lambda})_{\lambda \in L}\) . For every mapping \(\alpha: \mathbf{L} \to \mathbf{N}\) of finite support, we write

\[
e ^ {\alpha} = \prod_ {\lambda \in L} e _ {\lambda} ^ {\alpha (\lambda)}\tag{14}
\]

(product in the commutative algebra \(\mathbf{S}(\mathbf{M})\) ). Then, when \(\alpha\) runs through the set \(\mathbf{N}^{(L)}\) of mappings of \(\mathbf{L}\) into \(\mathbf{N}\) , of finite support, the \(e^{\alpha}\) form a basis of the A-module \(\mathbf{S}(\mathbf{M})\) .

As M is the direct sum of the \(M_{\lambda} = Ae_{\lambda}\) , it suffices to prove the theorem when L is reduced to a single element and then apply Proposition 9. But when M = Ae (e a free element), then \(x \otimes y - y \otimes x = 0\) for all x, y in M; the ideal \(J'\) is therefore zero, whence \(\mathsf{T}(\mathsf{A}e) = \mathsf{S}(\mathsf{A}e)\) and the theorem then follows from §5, no. 5, Theorem 1.

The multiplication table of the basis (14) is given by

\[
e ^ {\alpha} e ^ {\beta} = e ^ {\alpha + \beta}\tag{15}
\]

where \(\alpha + \beta\) is the mapping \(\lambda \mapsto \alpha(\lambda) + \beta(\lambda)\) of L into N. In other words, the basis \((e^{\alpha})\) of S(M), with the multiplicative law (15), is canonically isomorphic to the free commutative monoid \(\mathbf{N}^{(L)}\) derived from L; it follows (§ 2, no. 9) that the symmetric algebra S(M) of a free module M with basis whose indexing set is L, is canonically isomorphic to the polynomial algebra A{[}(X \(_{\lambda}\) ) \(_{\lambda \in L}\) {]}, the canonical isomorphism being obtained by mapping \(e_{\lambda}\) to X \(_{\lambda}\) . In particular (§ 2, no. 7, Proposition 7), for every mapping \(f\colon\) L \(\to\) E of L into a commutative A-algebra E, there exists one and only one A-algebra homomorphism \(\bar{f}\colon\) S(M) \(\to\) E such that \(\bar{f}(e_{\lambda}) = f(\lambda)\) .

Remark (2). The above results can equally well be obtained as a consequence of the universal properties of polynomial algebras and symmetric algebras, taking account of II, § 1, no. 11, Corollary 3 to Proposition 17.

COROLLARY. If M is a projective A-module, S(M) is a projective A-module.

The proof is the same as that for § 5, no. 5, Corollary to Theorem 1, replacing T by S.

PROPOSITION 10. Let \(\Delta\) be a commutative monoid, \(\mathbf{M}\) a graded A-module of type \(\Delta\) and \((\mathbf{M}_{\alpha})_{\alpha \in \Delta}\) its graduation. For every ordered pair \((\alpha, n) \in \Delta \times \mathbf{N}\) , let \(\mathbf{S}^{\alpha, n}(\mathbf{M})\) be the submodule of \(\mathbf{S}^n(\mathbf{M})\) the direct sum of the submodules \(\bigotimes_{\lambda \in J} \mathbf{S}^{n_{\lambda}}(\mathbf{M}_{\alpha_{\lambda}})\) , where \((n_{\lambda})_{\lambda \in L}\) runs through the set of families of integers \(\geqslant 0\) such that \(\sum_{\lambda \in L} n_{\lambda} = n, J\) is its support and, for each \((n_{\lambda})\) , \((\alpha_{\lambda})_{\lambda \in J}\) runs through the set of families of \(\Delta^J\) such that \(\sum_{\lambda \in J} \alpha_{\lambda} = \alpha\) . Then \((\mathbf{S}^{\alpha, n}(\mathbf{M}))_{(\alpha, n) \in \Delta \times \mathbf{N}}\) is the only graduation of type \(\Delta \times \mathbf{N}\) compatible with the algebra structure on \(\mathbf{S}(\mathbf{M})\) and which induces on \(\mathbf{M} = \mathbf{S}^1(\mathbf{M})\) the given graduation.

The fact that \(\mathbf{S}(\mathbf{M})\) is the direct sum of the \(\mathbf{S}^{\alpha,n}(\mathbf{M})\) follows from the Corollary to Proposition 9; the rest of the proof is identical with the end of the proof of §5, no. 5, Proposition 7.

Suppose more particularly that \(\Delta = \mathbf{Z}\) and let \(\mathbb{S}(\mathbf{M})\) be given the total graduation (of type \(\mathbf{Z}\) ) (II, §11, no. 1) corresponding to the graduation of type \(\mathbf{Z} \times \mathbf{N}\) (and hence also of type \(\mathbf{Z} \times \mathbf{Z}\) ) defined above; the homogeneous elements of degree \(n \in \mathbf{Z}\) under this graduation are therefore those of the direct sum of the \(\mathbb{S}^{q, m}(\mathbf{M})\) for \(q + m = n\) . Let \(f\) be a homogeneous linear mapping of degree 0 of the graded A-module \(\mathbf{M}\) into a commutative graded A-algebra \(\mathbf{F}\) of type \(\mathbf{Z}\) ; then the algebra homomorphism \(g: \mathbb{S}(\mathbf{M}) \to \mathbf{F}\) such that \(f = g \circ \phi_{\mathbf{M}}'\) is a homomorphism of graded algebras of type \(\mathbf{Z}\) , as follows from the formula \(g(x_1 x_2 \ldots x_n) = f(x_1) f(x_2) \ldots f(x_n)\) for homogeneous \(x_i\) in \(\mathbf{M}\) , from the hypothesis on \(f\) and from the definition of the graduation of type \(\mathbf{Z}\) on \(\mathbb{S}(\mathbf{M})\) .

